%% 04_agent_model.tex
\section{Agent Model}
\label{sec:agents}

This section gives the exact decision rule for each of the five
agent types in the simulator. All rules read the top of the book
(best bid $b_t$ and best ask $a_t$), the current fundamental
value~$\FV_t$, and a seeded pseudo-random draw $r\sim
\text{Uniform}(0,1)$; the formulas below are the literal rules
implemented in the source code, with constants exposed to the
reader so that any experiment can be reconstructed from this page
alone.

\subsection{Fundamentalist (F)}

The fundamentalist is a rational mean-reversion anchor inspired by
\citet{smith1988bubbles}. Its belief is
$\beta_t=\FV_t\cdot\bigl(1+0.02\,(r-0.5)\cdot 2\bigr)$---a
$\pm2\%$ symmetric noise about the fundamental---and its action
rule is
\begin{equation}
\label{eq:fund}
\begin{cases}
\text{cross the ask with a unit bid at }a_t, & \text{if }a_t<0.98\,\beta_t\text{ and cash}\ge a_t,\\
\text{cross the bid with a unit ask at }b_t, & \text{if }b_t>1.02\,\beta_t\text{ and inventory}>0,\\
\text{post passive with probability }0.45, & \text{otherwise,}\\
\text{hold}, & \text{otherwise.}
\end{cases}
\end{equation}
Passive quotes bid in $[0.970,\,1.005]\cdot\beta_t$ and ask in
$[0.995,\,1.030]\cdot\beta_t$, so the fundamentalist never posts a
cross on its own initiative. It supplies a narrow two-sided book
near the fundamental whenever neither side of the market is
aggressive and anchors mispricing back to $\FV_t$ through unit
crosses when either side drifts by more than~$2\%$.

\subsection{Trend follower (T)}

The trend follower estimates a short-horizon slope from the last
six non-null trade prices and projects three ticks forward,
\begin{equation}
\text{slope}_t\;=\;\frac{p_{t}-p_{t-h+1}}{\max(1,h-1)},\qquad
\text{target}_t\;=\;p_t+3\cdot\text{slope}_t,
\label{eq:trend}
\end{equation}
where $h=6$ is the slope window and $p_t$ is the most recent trade
price. When $\text{slope}_t>0$ and it has cash, the trend agent
chases up with a unit bid at $p_t\cdot(1+0.015+0.03\,r)$; when
$\text{slope}_t<0$ and it has inventory, it flees with a unit ask
at $p_t\cdot(1-0.015-0.03\,r)$. Two corner cases protect the agent
from stranding. First, if cash has dropped below the most recent
trade price but inventory still exceeds the initial endowment, the
agent dumps shares into the best bid at the aggressive price.
Second, in a flat-slope regime with excess inventory, the agent
liquidates with probability $\min(0.85,\,0.3+0.25\,e)$, where $e$
is the excess-inventory rate. The trend rule is the positive-feedback
channel \citep{delong1990positive} and is the primary amplifier of
bubble~formation in the simulator.

\subsection{Random zero-intelligence (R)}

The zero-intelligence trader posts uniform bids and asks on
fundamentally-scaled support, following
\citet{gode1993allocative}. Explicitly,
\begin{equation}
\text{bid}_t\sim\text{Uniform}[\,0.7\,\FV_t,\,1.1\,\FV_t\,],\qquad
\text{ask}_t\sim\text{Uniform}[\,0.9\,\FV_t,\,1.3\,\FV_t\,].
\label{eq:random}
\end{equation}
It never reads the order book, never tracks profit, and is pinned
in every run so that the book never goes fully empty. One slot of
this type is part of the fixed background and is invisible to the
population-composition sliders.

\subsection{Experienced (E)}

The experienced agent carries the DLM~2005 horizon discount. Let
$h=T-t+1$ be the number of periods remaining. Each experienced
agent carries an experience profile $p\in\{\text{naive},
\text{once-played},\text{veteran}\}$ with a corresponding pair
$(\varphi_0^{(p)},\varphi_1^{(p)})$ of a discount floor and a
discount slope, and its belief is
\begin{equation}
\beta_t^{(p)} \;=\; \FV_t\cdot\Bigl(\varphi_0^{(p)}+\frac{h}{T}\cdot\varphi_1^{(p)}\Bigr).
\label{eq:exp_belief}
\end{equation}
The three profiles are calibrated as in Table~\ref{tab:exp}.

\begin{table}[h]
\centering
\small
\caption{Experienced profiles}
\label{tab:exp}
\begin{tabular}{l c c c c}
\toprule
Profile & $\varphi_0$ & $\varphi_1$ & liquidate $h\le$ & avoid overpay $h\le$ \\
\midrule
Naive        & 1.00 & 0.05 & 1 & 2 \\
Once-played  & 0.96 & 0.04 & 2 & 3 \\
Veteran      & 0.92 & 0.06 & 3 & 4 \\
\bottomrule
\end{tabular}
\end{table}

When $h$ falls below the profile's liquidation threshold, the agent
dumps any remaining inventory into the best bid; when $h$ falls
below the overpay threshold, it refuses any buy above
$\beta_t^{(p)}$. Outside the terminal region it behaves like a
risk-tight fundamentalist, buying below $\beta_t^{(p)}$ with a
discount multiplier (\emph{e.g.\/}~0.92 for veterans) and selling
above $\beta_t^{(p)}$ with a premium multiplier (\emph{e.g.\/}~1.05).
The profile mix is user~adjustable, which allows the Experienced
population to be dialled between the DLM~2005 subject groups.

\subsection{Utility (U)}

The utility agent follows \citet{lopezlira2025llm}. Its prior at
every decision tick is
\begin{equation}
V_i^{\text{prior}}\;=\;\FV_t\cdot\bigl(1+b_i+\varepsilon_{i,t}\bigr),\qquad
b_i\in\{-0.15,\,0,\,+0.15\},\qquad
\varepsilon_{i,t}\sim\text{Uniform}[-0.03,\,0.03].
\label{eq:util_prior}
\end{equation}
The persistent bias $b_i$ is the ``persistently
optimistic/pessimistic/unbiased'' archetype and is assigned at
sampling time. The per-tick noise $\varepsilon$ keeps the agent's
quoting out of lockstep with other agents of the same~type.

The prior is then blended with incoming messages from peers. If
$\mathcal{M}_i$ denotes the set of claimed valuations received
during the current period, and $w\in(0,1)$ is a weight, the
posterior is
\begin{equation}
V_i^{\text{post}} \;=\; w\cdot V_i^{\text{prior}}
\;+\;(1-w)\cdot\frac{1}{|\mathcal{M}_i|}\sum_{m\in\mathcal{M}_i}\hat v_m.
\label{eq:blend}
\end{equation}
Three archetypes determine~$w$: \emph{naive} trusts peers and
uses $w=0.60$; \emph{skeptical} barely moves and uses $w=0.90$;
\emph{adaptive} scales $w$ by per-sender trust and caps the move at
$0.50$, with pairwise trust updated every trade by
\begin{equation}
\tau_{r\to s}\;\leftarrow\;(1-\alpha)\,\tau_{r\to s}
\;+\;\alpha\cdot\max\!\Bigl(0,\,1-\tfrac{|\hat v_s-\VWAP_t|}{\VWAP_t}\Bigr),\qquad
\alpha=0.30.
\label{eq:trust}
\end{equation}
Receiver~$r$'s trust in sender~$s$ is reinforced whenever $s$'s
claim is close to the period VWAP and attenuated~otherwise.

Given the posterior valuation $V_i^{\text{post}}$, the agent scores
each available action $a\in\{\text{hold},\text{cross bid},\text{cross ask},
\text{passive bid},\text{passive ask}\}$ by its expected utility
\begin{equation}
\EU(a)\;=\;p_{\text{fill}}\cdot U(w_1)\,+\,(1-p_{\text{fill}})\cdot U(w_0),
\qquad a^\star=\arg\max_{a}\EU(a),
\label{eq:eu}
\end{equation}
where $w_0$ and $w_1$ are the agent's wealth under non-fill and
fill respectively, and $p_{\text{fill}}=0.30$ is a constant passive
fill-probability inherited from the original Lopez-Lira design. The
risk-preference functional $U\in\{\Uloving,\Uneutral,\Uaverse\}$ is
assigned at sampling time and drawn from
\begin{equation}
\Uloving(w)=(w/w_0)^{2},\qquad
\Uneutral(w)=w/w_0,\qquad
\Uaverse(w)=\sqrt{w/w_0}.
\label{eq:risk}
\end{equation}

A deceptive variant of the utility rule is available for the
messaging channel. When enabled, an asset-poor agent (inventory
below its initial endowment) broadcasts a valuation equal to
$0.82\cdot V_i^{\text{true}}$, while an asset-rich agent broadcasts
$1.18\cdot V_i^{\text{true}}$. The lie-gap magnitude
$|\hat v^{\text{report}}-V^{\text{true}}|$ drives the period-end
trust update in~\eqref{eq:trust}, which produces the closed-loop
feedback described in \citet{lopezlira2025llm}.

Equations~\eqref{eq:fund}--\eqref{eq:risk} constitute the complete
decision model. The market aggregates these orders through a
price--time-priority order book, mutates cash and inventory on
every trade, credits dividends per \eqref{eq:dividend} at period
end, and records every state in append-only history arrays so the
replay system can reconstruct past ticks without re-running the
agents. No agent in this model consults a language model. The
next section adds the first such consultation as the Wang~2026
innovation.
